% Further E.3 deductions; owned by cycles. No preamble.
\paragraph{E.3g (Proved: integral abelian parameterization and regular invariants).}
Every algebraically trivial cycle $z\in\mathrm{CH}^i(X)_{\mathrm{alg}}$
on a smooth projective complex variety is, \emph{integrally}, the
fibre difference at a point and at zero of a cycle on $X\times A$
for some abelian variety $A$. Moreover every regular homomorphism
\[
 \rho:\mathrm{CH}^i(X)_{\mathrm{alg}}\longrightarrow B(\mathbb C)
\]
to an abelian variety kills $N^i(X)$.

For the first assertion, the definition of algebraic equivalence
expresses $z$ as a sum of $\Gamma_{j*}\delta_j$, with $C_j$ smooth
projective connected curves, $\delta_j\in\mathrm{CH}_0(C_j)^0$,
and $\Gamma_j\in\mathrm{CH}^i(C_j\times X)$. A normalized
Poincar\'e divisor on $J(C_j)\times C_j$ has fibre the divisor
class corresponding to any point of $J(C_j)=\operatorname{Pic}^0(C_j)$.
Composing this divisor correspondence with $\Gamma_j$ therefore
realizes $\Gamma_{j*}\delta_j$ exactly. Pulling the finitely many
families to $A=\prod_jJ(C_j)$ and adding them realizes $z$.
No isogeny denominator is involved. The up-to-an-integer argument
in E.3a remains necessary for its stronger assertion about
replacing a whole prescribed parameter variety by its own Albanese.

For the second assertion, regularity makes
$t\mapsto\rho(W_t-W_{t_0})$ a morphism from each smooth projective
connected parameter variety $T$ to $B$, taking $t_0$ to zero.
It therefore factors through $\operatorname{Alb}(T)$. Its value on
each $\nu\in\ker(\mathrm{CH}_0(T)^0\to\operatorname{Alb}(T))$
is zero. These images generate $N^i(X)$, proving the assertion.

The group $N^i(X)$ is divisible over $\mathbb C$. Indeed
$\mathrm{CH}_0(T)^0$ is divisible by curve parameterization. If
$\nu$ is in its Albanese kernel and $mw=\nu$, then
$\operatorname{alb}(w)\in\operatorname{Alb}(T)[m]$.
Roitman's torsion isomorphism
\citep[Theorem 4.2]{cyBloch1979} supplies
$\tau\in\mathrm{CH}_0(T)^0[m]$ with the same Albanese class.
Thus $w-\tau$ is in the Albanese kernel and $m(w-\tau)=\nu$.
Images and sums of divisible groups remain divisible, proving
the claim for $N^i(X)$. This is compatible with the independent
fact that the Walker kernel is divisible, whereas the ordinary
Abel--Jacobi kernel need not be
\citep[Theorem 2.3 and Remark 2.4, author p.~12]{cyACMVWalker2022}.

In particular the surjective regular Walker map
\[
 \psi_W:\mathrm{CH}^i(X)_{\mathrm{alg}}\longrightarrow J_W^i(X),
 \qquad
 J_W^i(X)=J\bigl(N^{i-1}H^{2i-1}(X,\mathbb Z(i))\bigr),
\]
factors through $P^i(X)$. Here $N^{i-1}$ in the cohomology group
is the coniveau filtration, not the subgroup $N^i(X)$ defining
Albanese equivalence. The regularity and surjectivity are the
opened \citet[Theorem 1.1, pp.~2--3; Corollaries 2.5--2.6,
pp.~7--8]{cySuzuki2021}. Thus we have
\[
 N^i(X)\subset\ker\psi_W\subset\ker\operatorname{AJ},
 \qquad
 \operatorname{AJ}=\alpha\psi_W,
\]
where $\alpha:J_W^i(X)\to J_a^i(X)$ is the natural isogeny.

\paragraph{E.3h (Refuted: integral equality with ordinary Abel--Jacobi equivalence).}
There exists a smooth projective complex fourfold $X$ and an
algebraically trivial codimension-three cycle $z$ such that
\[
 \operatorname{AJ}(z)=0,
 \qquad z\notin N^3(X).
\]
By E.3g this $z$ comes from an abelian family integrally. Hence the
integral equality of equivalence relations asserted on the right
side of Lead E4 is false. The same failure occurs for some
$d$-fold in every codimension $3\le i\le d-1$, $d\ge4$.

We use the exact opened input of
\citet[Corollary 4.2 and Remark 4.3, author p.~10]{cyOttemSuzuki2020}.
They construct a smooth threefold $W$, a rank-one degeneracy locus
of a map
\[
 \mathcal O^{\oplus3}\longrightarrow
 \mathcal O(1,2,0)\oplus\mathcal O(1,0,2)
 \quad\hbox{on }\mathbb P^1\times\mathbb P^2\times\mathbb P^2,
\]
and an elliptic curve $E$, for which $X=W\times E$ has a
surjective regular Walker map and the isogeny
$\alpha:J_W^3(X)\to J_a^3(X)$ has nonzero finite kernel.
More precisely, their construction takes coordinates $[S:T]$
on the first factor and the rank-one locus of the matrix
\[
 \begin{pmatrix}
 P_1&Q_1&R_1\\
 SP_2+qP_3&SQ_2+qQ_3&SR_2+qR_3
 \end{pmatrix}.
\]
Here $q$ is a sufficiently large auxiliary prime, and the three
triples of rational-coefficient polynomials have tridegrees
$(1,2,0)$, $(0,0,2)$ and $(1,0,2)$, respectively, with the
generality and specialization conditions imposed in their
Theorem 3.1. These are the exact construction data in that
theorem, not a selected numerical example.
Their Theorem 3.1 supplies $W$ over $\mathbb Q$ with torsion-free
integral cohomology and a nonalgebraic integral Hodge class; their
Theorem 4.1 supplies $E$. We use this as a proved primary theorem,
not as a claim that we have independently verified their full
degeneration argument or selected numerical coefficients for $W,E$.
E.3r below gives a fully displayed geometric generic matrix and
Legendre elliptic curve to which this obstruction transfers
integrally, and supplies a nonzero cycle of order two.

Choose $0\ne b\in\ker\alpha$ and use surjectivity to choose
$z$ with $\psi_W(z)=b$. Then $\operatorname{AJ}(z)=0$, while
E.3g gives $N^3(X)\subset\ker\psi_W$, so $z\notin N^3(X)$.
This is the complete deduction from the stated primary input.
For $d\ge4$ and $3\le i\le d-1$, put $s=d-4$, $j=i-3$ and
take $z'=z\times t^j$ on $X\times\mathbb P^s$. Its Abel--Jacobi
image is zero by functoriality. If $z'$ belonged to $N^i$, then
multiplication by $t^{s-j}$ followed by projection would put $z$
in $N^3(X)$, by functoriality of Albanese equivalence. Since
$0\le j\le s$, the projection is exactly $z$, a contradiction.

\paragraph{E.3i (Proved qualification: the finite obstruction does not settle
the family-map converse).}
The counterexample in E.3h does not itself give an abelian family
with $a({}^tZ)=0$ but $u_Z\ne0$. For any abelian family, regularity
gives an actual homomorphism
$\phi_{W,Z}:A\to J_W^i(X)$ with
$\alpha\phi_{W,Z}=\phi_Z$. If $a({}^tZ)=0$, then
$\phi_Z=0$ by the cylinder-map argument in E.3b. Thus the image
of $\phi_{W,Z}$ is contained in the finite group $\ker\alpha$.
An abelian variety is connected, so this homomorphism is zero.
Consequently the Walker invariant cannot detect a failure of the
family-map converse. The individual point $z$ in E.3h is realized
by a family whose whole Abel--Jacobi homomorphism need not vanish.

This identifies a precise logical defect in the unqualified
integral version of Lead E4: passage from a point in the kernel
of an abelian-variety homomorphism to its connected kernel requires
an integer multiple. E.3c proves the rational equivalence, and
E.3h refutes the proposed integral equality of equivalence relations.
We do not infer that the integral family-map converse is false.
E.3j--k below supply the additional argument needed to decide the
integral biconditional on this same fourfold.

\paragraph{E.3j (Proved: the family converse for one-cycles when zero-cycles
are supported on a curve).}
Let $X/\mathbb C$ be a smooth projective connected $n$-fold,
$n\ge2$, and suppose that a closed subvariety $C\subset X$
of dimension at most one supports $\mathrm{CH}_0(X)$; that is,
$\mathrm{CH}_0(C)\to\mathrm{CH}_0(X)$ is surjective. Then for
every abelian variety $A$ and every normalized
$Z\in\mathrm{CH}^{n-1}(A\times X)$,
\[
 a({}^tZ)=0\quad\Longrightarrow\quad u_Z=0
 \quad\text{in }\operatorname{Hom}\bigl(A(\mathbb C),P^{n-1}(X)\bigr).
\]
This is an integral conclusion about $u_Z$, although the
cohomological vanishing can be tested with rational coefficients.

\emph{The decomposition and its precise cycle action.}
We give the decomposition-of-the-diagonal argument needed here.
Choose a countable algebraically closed subfield $k_0\subset\mathbb C$
over which $X$ and $C$ are defined. Embed $K=k_0(X)$ in $\mathbb C$
over $k_0$. The corresponding geometric generic point $\eta$,
viewed as a complex point of $X$, is rationally equivalent to a
zero-cycle supported on $C$. The finitely many cycles and rational
functions in this relation spread over a finite-type integral
$\overline K$-scheme. A closed $\overline K$-point in a good open
specializes the relation to $\overline K$, preserving $\eta$ and
the support on $C$. The resulting finite data descend to a finite
extension $K'/K$. Taking the norm therefore gives
\[
 M[\eta]=z_C\quad\text{in }\mathrm{CH}_0(X_K),
 \qquad M>0,\quad\operatorname{supp}(z_C)\subset C_K.
\]
Closing $z_C$ in $X\times C$ and applying localization to the
generic point of the first factor yields
\begin{equation}\label{eq:E3-support-decomposition}
 M\Delta_X=\Gamma_C+\Gamma_D\quad\text{in }\mathrm{CH}^n(X\times X),
\end{equation}
where $\Gamma_C$ is supported on $X\times C$ and $\Gamma_D$ on
$D\times X$ for a divisor $D\subset X$. We have proved the
needed decomposition from the stated complex support hypothesis;
we have not assumed universal support over every field extension.

Normalize the curve components of $C$. Their contribution lifts
to divisor correspondences on $X\times\widetilde C$; components
over isolated points have the same, or simpler, factorization.
They act on $\mathrm{CH}^{n-1}(X)$ through
$\mathrm{CH}^{0}(\widetilde C)$, since the lifted correspondence
has codimension one and the source has dimension $n$.
As $\mathrm{CH}^{0}(\widetilde C)$ has no algebraically trivial
part, $\Gamma_C$ annihilates
$\mathrm{CH}^{n-1}(X)_{\mathrm{alg}}$.

For each component of $D$ choose a projective resolution
$f_j:T_j\to D_j\hookrightarrow X$, with $T_j$ smooth of
dimension $n-1$. After rational coefficients, every cycle
supported on $D\times X$ lifts to a sum of cycles on
$T_j\times X$. To check surjectivity even over singular strata,
take a component of the inverse image of the given cycle that
dominates it, then a projective multisection of its generic
fibre. The closure is generically finite of positive degree;
division by that degree gives the desired lift. Clearing the
finitely many denominators in \eqref{eq:E3-support-decomposition},
we obtain an integer, still denoted $M>0$, and integral cycles
$G_j\in\mathrm{CH}^{n-1}(T_j\times X)$ such that, on
algebraically trivial cycles,
\[
 M\operatorname{id}
       =\sum_j(G_j)_*f_j^*.
\]
All these operations preserve Albanese equivalence by E.3a.
Thus the same identity factors through
\begin{equation}\label{eq:E3-Alb-factorization}
 P^{n-1}(X)\xrightarrow{(f_j^*)_j}
       \prod_j\operatorname{Alb}(T_j)(\mathbb C)
       \xrightarrow{\sum_j(G_j)_*}P^{n-1}(X).
\end{equation}
The composition is multiplication by $M$; it is not asserted
to be the identity integrally.

\emph{Apply this factorization to the family.}
The pullback $(\operatorname{id}_A\times f_j)^*Z$ is a
zero-cycle family on $T_j$. If $a({}^tZ)=0$, its cylinder map
is zero too, by functoriality: it is the original cylinder
map followed by
$f_j^*:H^{2n-3}(X,\mathbb Q)\to H^{2n-3}(T_j,\mathbb Q)$.
The corresponding homomorphism
$A\to\operatorname{Alb}(T_j)$ is therefore zero by faithfulness
on $H_1$, or by the zero-cycle endpoint case E.3d.
Equation \eqref{eq:E3-Alb-factorization} gives $M u_Z=0$.
Finally $A(\mathbb C)$ is divisible: for every $a$ there is
$b$ with $Mb=a$. Hence $u_Z(a)=M u_Z(b)=0$. This proves
the integral family-map converse under the stated support
hypothesis.

\paragraph{E.3k (Refuted: the unqualified integral biconditional in Lead E4;
proved corrected fixed-target criterion).}
For the fourfold $X=W\times E$ of E.3h, the map
$\mathrm{CH}_0(W)\to\mathbb Z$ is an isomorphism by
\citet[Lemma 2.4, author p.~4]{cyOttemSuzuki2020}.
Fix $w_0\in W(\mathbb C)$. For every point $(w,e)$,
the equality $[w]=[w_0]$ on $W$, pushed into $W\times\{e\}$,
shows that its zero-cycle class is represented on
$\{w_0\}\times E$. Thus this curve supports
$\mathrm{CH}_0(X)$, and E.3j proves the family-map converse
for every codimension-three abelian family on this $X$.
E.3h simultaneously gives an algebraically trivial cycle from
an abelian family that is Abel--Jacobi trivial but is not
Albanese trivial. The two assertions declared equivalent in
the literal integral Lead E4 therefore have different truth
values on the same smooth projective fourfold. This refutes
that biconditional, while leaving E.3c's rational formulation
intact.

For an arbitrary smooth projective complex $X$, a precise
fixed-target criterion uses
\[
 K_W^i(X)=\ker\bigl(P^i(X)\longrightarrow J_W^i(X)(\mathbb C)\bigr).
\]
The integral family-map converse is exactly the assertion
that no nonzero image $u_Z(A(\mathbb C))$ of an abelian family
is contained in $K_W^i(X)$. Indeed such containment is
equivalent to $\phi_{W,Z}=0$, which is equivalent to
$a({}^tZ)=0$: one direction uses the finite-kernel argument
in E.3i, and the other the cylinder map in E.3b.
The rational criterion in E.3c is equivalently
$K_W^i(X)\otimes\mathbb Q=0$, because
\[
 0\longrightarrow K_W^i(X)\longrightarrow
 \ker\overline{\operatorname{AJ}}
 \longrightarrow\ker\alpha\longrightarrow0
\]
is exact and $\ker\alpha$ is finite; surjectivity on the right
follows from surjectivity of $\psi_W$.
Neither condition follows merely from regularity or from the
existence of a universal regular homomorphism. We do not
identify $P^i(X)$ with the Walker Jacobian in general.

\paragraph{E.3l (Proved: the integral Walker quotient for one-cycles
under support on a surface).}
Let $X/\mathbb C$ be a smooth projective connected $n$-fold,
$n\ge2$, and assume that $\mathrm{CH}_0(X)$ is supported on
a closed subvariety of dimension at most two. Then the
Walker map induces a canonical isomorphism of abelian groups
\begin{equation}\label{eq:E3-Walker-quotient}
 P^{n-1}(X)\ \xrightarrow{\ \sim\ }\ 
               J_W^{n-1}(X)(\mathbb C).
\end{equation}
In particular the integral family-map converse holds in this
codimension, and its discrepancy from ordinary Abel--Jacobi
equivalence is exactly the finite kernel of
$\alpha:J_W^{n-1}(X)\to J_a^{n-1}(X)$.

\emph{Proof.}
The generic-point, norm, and localization proof in E.3j
applies verbatim with a surface $S$ in place of $C$.
Resolve its two-dimensional components, and lift the
correspondence supported on $X\times S$ by the same
multisection argument. A lift to $X\times\widetilde S$ has
codimension two, so on codimension-$n-1$ algebraically trivial
cycles it acts through
\[
 \mathrm{CH}^1(\widetilde S)_{\mathrm{alg}}
              =\operatorname{Pic}^0(\widetilde S)(\mathbb C).
\]
Lower-dimensional components act through
$\mathrm{CH}^0$ of curves or zero groups and therefore
annihilate algebraically trivial cycles. The term supported
on $D\times X$ still factors through the Albanese varieties
of the resolutions $T_j$ of $D$.

More explicitly, write $j_a:\widetilde S_a\to X$ for the
resolved surface maps. After clearing denominators we can
write the surviving terms of the diagonal as
$(\operatorname{id}_X\times j_a)_*H_a$ and
$(f_j\times\operatorname{id}_X)_*G_j$, where
\[
 H_a\in\mathrm{CH}^2(X\times\widetilde S_a),\qquad
 G_j\in\mathrm{CH}^{n-1}(T_j\times X).
\]
The projection formula, or composition with the respective
graph correspondences, identifies their actions exactly as
$(j_a)_*(H_a)_*$ and $(G_j)_*f_j^*$. Here $f_j^*$ is the
pullback in the Chow ring between smooth varieties, defined
by the graph's refined Gysin map; it is defined even though
$D$ need not be smooth. In particular
\[
\begin{aligned}
 (H_a)_*&: \mathrm{CH}^{n-1}(X)_{\mathrm{alg}}
                \longrightarrow\operatorname{Pic}^0(\widetilde S_a),\\
 f_j^*&: \mathrm{CH}^{n-1}(X)_{\mathrm{alg}}
                \longrightarrow\mathrm{CH}_0(T_j)^0.
\end{aligned}
\]
Composing the second map with $\operatorname{alb}_{T_j}$
therefore gives the asserted factors. Every correspondence
preserves $N$: compose it with each generator's correspondence
in the definition of $N$, leaving the same Albanese-kernel
zero-cycle on its parameter variety. At the two endpoints,
$P^1(\widetilde S_a)=\operatorname{Pic}^0(\widetilde S_a)$ and
$P^{n-1}(T_j)=\operatorname{Alb}(T_j)(\mathbb C)$ by E.3d.
Thus both forward and return maps below are well-defined on
the indicated quotients. After clearing denominators,
we have a factorization of multiplication by an integer $M>0$,
\[
 B=\prod_a\operatorname{Pic}^0(\widetilde S_a)
       \times\prod_j\operatorname{Alb}(T_j),
 \qquad
 P^{n-1}(X)\xrightarrow{F}B(\mathbb C)
 \xrightarrow{G}P^{n-1}(X),
\]
with $GF=M\operatorname{id}$.
Choose an algebraically trivial representative $z$ of a class
in $\ker\overline{\operatorname{AJ}}$. Ordinary Abel--Jacobi
functoriality for correspondences gives, integrally,
\[
\begin{aligned}
 \operatorname{AJ}^1_{\widetilde S_a}((H_a)_*z)
       &=(H_a)_{\operatorname{AJ}}\operatorname{AJ}^{n-1}_X(z)=0,\\
 \operatorname{alb}_{T_j}(f_j^*z)
       &=(f_j^*)_{\operatorname{AJ}}\operatorname{AJ}^{n-1}_X(z)=0.
\end{aligned}
\]
The first Abel--Jacobi map is the identification of
algebraically trivial divisors with $\operatorname{Pic}^0$;
the second is the zero-cycle Abel--Jacobi map, namely the
Albanese map. Hence every component of $F(z)$ is zero.
This uses ordinary Abel--Jacobi functoriality, and does not
require an integral inverse to an isogeny. Consequently
\begin{equation}\label{eq:E3-bounded-AJ-kernel}
             M\ker\overline{\operatorname{AJ}}=0.
\end{equation}

On the other hand, the exact primary statement
\citet[Theorem 2.3, author p.~12]{cyACMVWalker2022} says
that the \emph{whole} kernel of $\psi_W$ is divisible.
Since $N^{n-1}(X)\subset\ker\psi_W$, the group
\[
 K_W^{n-1}(X)=\ker\psi_W/N^{n-1}(X)
\]
is divisible. It is a subgroup of
$\ker\overline{\operatorname{AJ}}$, so
\eqref{eq:E3-bounded-AJ-kernel} kills it by $M$.
A divisible group killed by a positive integer is zero:
write each element as $M$ times another element. Thus
$K_W^{n-1}(X)=0$. Surjectivity of $\psi_W$ proves
\eqref{eq:E3-Walker-quotient} integrally. The exact sequence
in E.3k now identifies the ordinary Abel--Jacobi kernel on
$P^{n-1}(X)$ with $\ker\alpha$ itself.

Under this isomorphism, a family map $u_Z$ is precisely the
regular homomorphism $\phi_{W,Z}$ on complex points. If its
cohomological action is zero, E.3i proves
$\phi_{W,Z}=0$, so $u_Z=0$. This also proves the converse
for the full support-on-a-surface class. It does not extend
the quotient identification to intermediate codimensions
without the stated support hypothesis.
